% !TeX root = ../../main.tex
% ---------------------------------------------------------------------------
% Deterministic geometric violet floral page borders (presentation layer).
%
% Design contract:
%   * Vector-only TikZ ornament drawn on a shipout background layer (eso-pic),
%     so it never affects bounding boxes, line breaking, or pagination.
%   * All ornament is hard-clipped to two vertical margin bands, one along
%     each page edge: x in [FloralBandMin, FloralBandMax] mm from the edge,
%     y in [FloralBandVPad, paperheight-FloralBandVPad] mm.  With letterpaper
%     and 0.82in (20.8mm) text margins this leaves >= 6.8mm of empty space
%     between any ornament and the content rectangle, keeps the running
%     header (which starts 20.8mm from the left edge) untouched, and leaves
%     the bottom-center page-number zone completely empty.
%   * Determinism: page_seed = FloralBaseSeed + FloralSeedPrime * abspage,
%     where abspage is a private absolute shipout counter (independent of
%     roman/arabic numbering).  The PGF random seed is reset from this value
%     at the start of every physical page, so no page's art depends on
%     earlier drawing.  BASE = 104729, PRIME = 9973 (both fixed; documented
%     here as required).  Randomness selects only among prevalidated slots
%     and bounded parameter ranges -- never free coordinates.
%   * Disable everything with \floralbordersfalse (before \begin{document}).
% ---------------------------------------------------------------------------

\newif\iffloralborders
\floralborderstrue
\makeatletter

% --- fixed seeds (see header comment) --------------------------------------
\def\FloralBaseSeed{104729}
\def\FloralSeedPrime{9973}

% --- master knobs -----------------------------------------------------------
\def\FloralOpacityScale{1.0}   % master opacity multiplier
\def\FloralFrameX{7.0}         % frame-line inset from page edge (mm)
\def\FloralBandMin{2.0}        % clip band, min distance from edge (mm)
\def\FloralBandMax{14.0}       % clip band, max distance from edge (mm)
\def\FloralBandVPad{7.0}       % clip band, vertical padding top/bottom (mm)
\def\FloralMotifX{9.4}         % nominal motif-center distance from edge (mm)

% --- palette (defined once; violet family + one restrained gold accent) -----
\definecolor{floralDeep}{HTML}{4A2E63}   % deep violet/plum: outlines, centers
\definecolor{floralMid}{HTML}{7C4E9F}    % medium violet: primary petals
\definecolor{floralLav}{HTML}{A88FC7}    % muted lavender: secondary petals
\definecolor{floralLilac}{HTML}{D9CBE9}  % pale lilac: low-opacity accents
\definecolor{floralGold}{HTML}{A9853B}   % rare complementary accent

% --- absolute page counter ---------------------------------------------------
\newcount\floral@abspage
\newcount\floral@seed
\newcount\floral@seedtmp

% --- bounded random helpers (pgf stream, reseeded per page) ------------------
% \floral@rnd{\macro}{min}{max} -> integer in [min,max]
\def\floral@rnd#1#2#3{\pgfmathrandominteger{#1}{#2}{#3}}

% ============================================================================
% Motif families (all built from circles, regular polygons, arcs, segments).
% Each draws around origin at unit scale ~= 1mm, max radius ~2.8.
% Argument #1: palette variant 0/1 (primary violet vs lavender).
% ============================================================================

% Family A: radial circle-petal blossom -- six petal circles around a center.
\newcommand\floral@motifA[1]{%
  \ifnum#1=0\colorlet{fpet}{floralMid}\else\colorlet{fpet}{floralLav}\fi
  \foreach \a in {0,60,...,300}{%
    \fill[fpet,opacity=0.38*\FloralOpacityScale] (\a:1.75) circle (0.85);
    \draw[floralDeep,opacity=0.45*\FloralOpacityScale,line width=0.35pt]
      (\a:1.75) circle (0.85);}%
  \fill[floralDeep,opacity=0.62*\FloralOpacityScale] (0,0) circle (0.62);}

% Family B: faceted hexagonal blossom -- two nested hexagons, offset 30deg.
\newcommand\floral@motifB[1]{%
  \ifnum#1=0\colorlet{fpet}{floralMid}\else\colorlet{fpet}{floralLav}\fi
  \draw[fpet,opacity=0.55*\FloralOpacityScale,line width=0.5pt]
    (0:2.4)--(60:2.4)--(120:2.4)--(180:2.4)--(240:2.4)--(300:2.4)--cycle;
  \fill[floralLilac,opacity=0.40*\FloralOpacityScale]
    (30:1.45)--(90:1.45)--(150:1.45)--(210:1.45)--(270:1.45)--(330:1.45)--cycle;
  \draw[floralDeep,opacity=0.50*\FloralOpacityScale,line width=0.35pt]
    (30:1.45)--(90:1.45)--(150:1.45)--(210:1.45)--(270:1.45)--(330:1.45)--cycle;
  \fill[floralDeep,opacity=0.62*\FloralOpacityScale] (0,0) circle (0.42);}

% Family C: offset polygon rosette -- five rotated diamonds around a center.
\newcommand\floral@motifC[1]{%
  \ifnum#1=0\colorlet{fpet}{floralMid}\else\colorlet{fpet}{floralLav}\fi
  \foreach \a in {90,162,234,306,18}{%
    \fill[fpet,opacity=0.34*\FloralOpacityScale,rotate=\a]
      (1.0,0)--(1.9,0.55)--(2.8,0)--(1.9,-0.55)--cycle;
    \draw[floralDeep,opacity=0.42*\FloralOpacityScale,line width=0.3pt,rotate=\a]
      (1.0,0)--(1.9,0.55)--(2.8,0)--(1.9,-0.55)--cycle;}%
  \draw[floralDeep,opacity=0.60*\FloralOpacityScale,line width=0.45pt]
    (0,0) circle (0.55);}

% Family D: arc blossom -- five open petal arcs; reads well when edge-clipped.
\newcommand\floral@motifD[1]{%
  \ifnum#1=0\colorlet{fpet}{floralMid}\else\colorlet{fpet}{floralLav}\fi
  \foreach \a in {90,162,234,306,18}{%
    \draw[fpet,opacity=0.55*\FloralOpacityScale,line width=0.55pt,rotate=\a]
      (0.7,0.55) arc[start angle=140,end angle=40,radius=1.35];}%
  \fill[floralDeep,opacity=0.58*\FloralOpacityScale] (0,0) circle (0.5);
  \draw[floralLav,opacity=0.45*\FloralOpacityScale,line width=0.35pt]
    (0,0) circle (0.95);}

% Family E: stem-and-bud constellation -- curved stem with bud nodes.
\newcommand\floral@motifE[1]{%
  \ifnum#1=0\colorlet{fpet}{floralMid}\else\colorlet{fpet}{floralLav}\fi
  \draw[floralDeep,opacity=0.48*\FloralOpacityScale,line width=0.45pt]
    (0,-2.6) .. controls (0.55,-1.0) and (-0.45,0.6) .. (0.1,2.4);
  \fill[fpet,opacity=0.55*\FloralOpacityScale] (0.1,2.4) circle (0.55);
  \fill[floralLilac,opacity=0.55*\FloralOpacityScale] (0.42,-0.35) circle (0.34);
  \fill[fpet,opacity=0.42*\FloralOpacityScale,rotate around={35:(-0.35,-1.4)}]
    (-0.35,-1.4) ++(-0.3,0) -- ++(0.3,0.42) -- ++(0.3,-0.42) -- ++(-0.3,-0.42)
    -- cycle;}

% Family F: three-node floral cluster -- tiny circle/diamond/triangle linked.
\newcommand\floral@motifF[1]{%
  \ifnum#1=0\colorlet{fpet}{floralMid}\else\colorlet{fpet}{floralLav}\fi
  \draw[floralLav,opacity=0.42*\FloralOpacityScale,line width=0.35pt]
    (-1.1,-1.5) -- (0.2,0.1) -- (0.9,1.8);
  \fill[fpet,opacity=0.55*\FloralOpacityScale] (0.9,1.8) circle (0.5);
  \fill[floralDeep,opacity=0.55*\FloralOpacityScale]
    (0.2,0.1) ++(0,0.42) -- ++(0.38,-0.42) -- ++(-0.38,-0.42) -- ++(-0.38,0.42)
    -- cycle;
  \fill[floralLilac,opacity=0.62*\FloralOpacityScale]
    (-1.1,-1.5) ++(0,0.42) -- ++(0.38,-0.65) -- ++(-0.76,0) -- cycle;}

% Dispatch: draw family #1 (0..5), palette variant #2, at (#3,#4) mm,
% scale #5 (per-mille -> /1000 later is avoided: #5 already decimal), rot #6.
\newcommand\floral@motif[6]{%
  \begin{scope}[shift={(#3,#4)},scale=#5,rotate=#6]
    \ifcase#1\floral@motifA{#2}\or\floral@motifB{#2}\or\floral@motifC{#2}\or
      \floral@motifD{#2}\or\floral@motifE{#2}\or\floral@motifF{#2}\fi
  \end{scope}}

% ============================================================================
% Per-page composition
% ============================================================================
\newcommand\floral@drawpage{%
  \begin{tikzpicture}[overlay,x=1mm,y=1mm]
    \pgfmathsetmacro\fpw{\paperwidth/1mm}%
    \pgfmathsetmacro\fph{\paperheight/1mm}%
    % ---- reseed deterministically from the absolute page ----
    % page_seed = (BASE + PRIME*abspage) mod 16000  (pgfmath seed range limit;
    % \divide truncates, so this is a true modulo)
    \floral@seed=\numexpr\FloralBaseSeed+\FloralSeedPrime*\floral@abspage\relax
    \floral@seedtmp=\floral@seed
    \divide\floral@seedtmp by 16000
    \multiply\floral@seedtmp by 16000
    \advance\floral@seed by -\floral@seedtmp
    \pgfmathsetseed{\the\floral@seed}%
    % ---- hard clip: union of the two vertical edge bands ----
    \clip (\FloralBandMin,\FloralBandVPad)
            rectangle (\FloralBandMax,\fph-\FloralBandVPad)
          (\fpw-\FloralBandMax,\FloralBandVPad)
            rectangle (\fpw-\FloralBandMin,\fph-\FloralBandVPad);
    % ---- page class: 2=title, 1=normal, 0=quiet (TOC 3-4, bibliography 32-34)
    \def\floral@class{1}%
    \ifnum\floral@abspage=1 \def\floral@class{2}\fi
    \ifnum\floral@abspage=3 \def\floral@class{0}\fi
    \ifnum\floral@abspage=4 \def\floral@class{0}\fi
    \ifnum\floral@abspage>31 \def\floral@class{0}\fi
    % ---- dominant (busier) side: outer edge alternates with page parity ----
    \pgfmathsetmacro\floral@odd{mod(\floral@abspage,2)}%
    % dominx = motif-center x on the dominant side; quiet side mirrored
    % ---- segmented frame lines (thin, lighter than flower centers) ----
    \floral@rnd\fgap{55}{75}%   gap position as % of usable height
    \pgfmathsetmacro\fylo{18}\pgfmathsetmacro\fyhi{\fph-18}%
    \pgfmathsetmacro\fgapy{\fylo+(\fyhi-\fylo)*(\fgap/100)}%
    \ifdim\floral@odd pt>0.5pt
      \pgfmathsetmacro\fdomx{\fpw-\FloralFrameX}\pgfmathsetmacro\fquix{\FloralFrameX}%
    \else
      \pgfmathsetmacro\fdomx{\FloralFrameX}\pgfmathsetmacro\fquix{\fpw-\FloralFrameX}%
    \fi
    \draw[floralLav,opacity=0.38*\FloralOpacityScale,line width=0.45pt]
      (\fdomx,\fylo) -- (\fdomx,\fgapy-9);
    \draw[floralLav,opacity=0.38*\FloralOpacityScale,line width=0.45pt]
      (\fdomx,\fgapy+9) -- (\fdomx,\fyhi);
    \draw[floralLilac,opacity=0.45*\FloralOpacityScale,line width=0.4pt]
      (\fquix,{0.36*\fph}) -- (\fquix,{0.64*\fph});
    % small terminal marks on the dominant frame line
    \fill[floralMid,opacity=0.5*\FloralOpacityScale] (\fdomx,\fylo) circle (0.5);
    \fill[floralMid,opacity=0.5*\FloralOpacityScale] (\fdomx,\fyhi) circle (0.5);
    % ---- motif slots: five prevalidated heights per side ----
    % slot y fractions: 12, 30, 50, 70, 88 percent of page height
    \def\floral@sloty##1{%
      \ifcase##1 0.12\or0.30\or0.50\or0.70\or0.88\fi}%
    % dominant slot index rotates deterministically so consecutive pages
    % never repeat the same dominant corner/height
    \pgfmathtruncatemacro\fdomslot{mod(\floral@abspage*2,5)}%
    % motif count by class
    \ifnum\floral@class=2 \def\fmotifn{4}\else
      \ifnum\floral@class=0 \def\fmotifn{1}\else
        \floral@rnd\fmotifn{2}{3}\fi\fi
    % draw motifs
    \foreach \fk in {1,...,\fmotifn}{%
      \floral@rnd\ffam{0}{5}%
      \floral@rnd\fpal{0}{1}%
      \floral@rnd\frot{0}{359}%
      \floral@rnd\fscl{75}{115}%
      \floral@rnd\fjy{-40}{40}%   y jitter, tenths of mm
      \floral@rnd\fjx{-10}{10}%   x jitter, tenths of mm
      \pgfmathsetmacro\fsc{\fscl/100}%
      \ifnum\fk=1
        % dominant motif: deterministic slot on the dominant side
        \pgfmathtruncatemacro\fslot{\fdomslot}%
        \ifdim\floral@odd pt>0.5pt
          \pgfmathsetmacro\fmx{\fpw-\FloralMotifX+\fjx/10}%
        \else
          \pgfmathsetmacro\fmx{\FloralMotifX+\fjx/10}%
        \fi
      \else
        \floral@rnd\fslot{0}{4}%
        \floral@rnd\fside{0}{1}%
        \ifnum\fside=0
          \pgfmathsetmacro\fmx{\FloralMotifX+\fjx/10}%
        \else
          \pgfmathsetmacro\fmx{\fpw-\FloralMotifX+\fjx/10}%
        \fi
      \fi
      \pgfmathsetmacro\fmy{\floral@sloty{\fslot}*\fph+\fjy/10}%
      \floral@motif{\ffam}{\fpal}{\fmx}{\fmy}{\fsc}{\frot}%
    }%
    % ---- occasional intentionally edge-clipped rosette (normal pages) ----
    \ifnum\floral@class=1
      \floral@rnd\fclipq{1}{4}%
      \ifnum\fclipq=1
        \floral@rnd\fcy{20}{80}%
        \floral@rnd\fcs{0}{1}%
        \ifnum\fcs=0 \pgfmathsetmacro\fcx{1.2}\else
          \pgfmathsetmacro\fcx{\fpw-1.2}\fi
        \floral@motif{3}{1}{\fcx}{\fcy/100*\fph}{1.05}{0}%
      \fi
    \fi
    % ---- tiny accent nodes (pale lilac; gold is rare) ----
    \ifnum\floral@class=0 \def\faccn{1}\else\floral@rnd\faccn{1}{3}\fi
    \foreach \fk in {1,...,\faccn}{%
      \floral@rnd\fay{15}{85}%
      \floral@rnd\fas{0}{1}%
      \floral@rnd\fgold{1}{12}%
      \ifnum\fas=0 \pgfmathsetmacro\fax{\FloralFrameX}\else
        \pgfmathsetmacro\fax{\fpw-\FloralFrameX}\fi
      \ifnum\fgold=1
        \fill[floralGold,opacity=0.55*\FloralOpacityScale]
          (\fax,\fay/100*\fph) circle (0.42);
      \else
        \fill[floralLilac,opacity=0.62*\FloralOpacityScale]
          (\fax,\fay/100*\fph) circle (0.42);
      \fi}%
    % ---- title page: two extra corner blossoms + gold buds ----
    \ifnum\floral@class=2
      \floral@motif{0}{0}{\FloralMotifX}{0.88*\fph}{1.15}{15}%
      \floral@motif{1}{1}{\fpw-\FloralMotifX}{0.12*\fph}{1.15}{0}%
      \fill[floralGold,opacity=0.55*\FloralOpacityScale]
        (\FloralFrameX,{0.80*\fph}) circle (0.45);
      \fill[floralGold,opacity=0.55*\FloralOpacityScale]
        (\fpw-\FloralFrameX,{0.20*\fph}) circle (0.45);
    \fi
  \end{tikzpicture}}

% ---- hook the drawing into the shipout background layer ---------------------
\AddToShipoutPictureBG{%
  \iffloralborders
    \global\advance\floral@abspage by 1
    \AtPageLowerLeft{\floral@drawpage}%
  \fi}
\makeatother
